% Week 13. LEC08 PDF pp. 17--31, transcript 79752; pp. 32--43, transcript 79753.
\section{第 13 节课：立体匹配与三维重建}
\label{lecture:SC4061-L13}

\lecturemeta
  {SC4061-L13}
  {2026-11-10}
  {LEC08 Stereo Vision -- 2026}
  {pp.~17--31（跨视图匹配）；pp.~32--43（三维重建）}
  {字幕 79752（对应点）；79753（三角测量与重建实例）}
  {已核对讲义与字幕}

上一讲用 $Z=fT/d$ 将视差换算成深度。这一步有个前提：量视差的两个像素必须来自同一个空间点。
本讲先解决怎样在右图找到这个点，再讨论两台相机朝向不同时怎样恢复它的三维坐标。

\lecturetopic{SC4061-L13-T01}{沿同一行寻找对应点：外观匹配与 SSD}

经过\relatedtopic{SC4061-L12-T04}{水平立体校正}，左图像素 $(u,v)$ 在右图的对应点也在第 $v$ 行，横坐标仍然未知。
设左右图的主点已对齐，用 $d=u_L-u_R$ 表示候选视差。取左图的一小块作为模板，在右图同一行上移动窗口，逐像素相减、平方再求和，得到平方差之和（SSD）：
\[
 C(d)=\sum_{(i,j)\in\Omega}
 \big[I_L(u+i,v+j)-I_R(u-d+i,v+j)\big]^2,
 \qquad \hat d\in\arg\min_{d\in D} C(d).
\]
这里 $\Omega$ 是窗口内的像素偏移，$D$ 是允许搜索且窗口有效的视差集合。
讲义 p.~21 给出的是对候选中心 $(x,y)$ 的一般搜索；水平校正把它缩成沿一行搜索。
把每个能可靠匹配的位置所测得的 $\hat d$ 放回左图，就得到视差图。SSD 的手算见\relatedexercise{SC4061-L13-E01}{两个候选窗口的比较}。

\textbf{哪些位置能匹配？}如果两图都拍到了同一表面，局部形变小，照明稳定，而且表面近似漫反射，两个窗口往往很像。能否从中挑出一个位置，还取决于纹理：
\begin{itemize}
  \item 一面均匀的白墙上，沿行移动窗口，许多位置的 SSD 都接近零；最低分也不足以确定视差。
  \item 重复的窗格会产生多个相似的低谷。带有独特花纹的窗口更容易得到一个清楚的最低点。
  \item 若右相机被前景挡住，它根本没有拍到该表面点。即使算法能选出最低分，把这一对送去测深也没有几何依据。
\end{itemize}

\textbf{基线怎样选？}较短的基线 $T$ 往往让左右外观更相似；拉开相机能增加视差，但也可能加重遮挡和外观变化。
沿用上一讲的小误差近似，$|\delta Z|\approx Z^2|\delta d|/(fT)$。
对同一深度、同一焦距的点，若基线加倍且像素视差误差也加倍，分子和分母中的倍数抵消，一阶绝对深度误差保持不变。

\lecturetopic{SC4061-L13-T02}{视点变化较大时：匹配局部特征}

当同一物体从侧面看发生明显形变时，逐像素比较的代价会上升。特征匹配先在两图中各自找出角点、曲线上的特征点等，再比较这些点附近的局部形状或描述向量。这样得到的是一组带有左右图像坐标的\textbf{稀疏对应点}。

例如，一个角点在图上旋转之后，角的两条边朝向都变了，夹角却可保持相近。讲义还用曲线的局部曲率作例子。视点变化过大时，这些量也可能改变。
\relatedtopic{SC4061-L09-T03}{SIFT 描述子}则把关键点附近的梯度按局部方向对齐，在 $4\times4$ 个小格中各统计 8 个方向，组成 128 维向量。左右图的描述子越相似，相应关键点越值得作为候选。

若同一轮廓上有好几个相似的角，可以参考两条经验规则：优先考虑在两图中位置相近的候选（proximity），或检查一串特征沿轮廓的先后次序能否对应（ordering）。遮挡会让部分特征消失，近处物体也可能发生很大的像素位移，因此这两条规则都需要结合图像判断。见\relatedexercise{SC4061-L13-E02}{可见性与匹配方法}。

\lecturetopic{SC4061-L13-T03}{从两幅图的对应点求三维坐标}

已知两台相机在同一世界坐标系下的投影矩阵，一个像素给出从相机光心出发的一条视线。理想情况下，同一空间点在两图中对应的视线相交，交点就是要找的 $\bm X=(X,Y,Z)^T$。这套方法允许两台相机有不同的朝向；若标定时采用米作长度单位，恢复的坐标也以米计。

\textbf{每幅图贡献两个方程。}把一台相机已知的 $3\times4$ 投影矩阵 $P$ 的第 $i$ 行写成 $[\,\bm r_i^T\mid b_i\,]$。按讲义 pp.~34--37 的投影关系展开，有
\[
 u=\frac{\bm r_1^T\bm X+b_1}{\bm r_3^T\bm X+b_3},\qquad
 v=\frac{\bm r_2^T\bm X+b_2}{\bm r_3^T\bm X+b_3}.
\]
对于有限像点，把分母乘过去，再将含 $\bm X$ 的项移到左边：
\[
 (\bm r_1-u\bm r_3)^T\bm X=u b_3-b_1,\qquad
 (\bm r_2-v\bm r_3)^T\bm X=v b_3-b_2.
\]
左右相机各代入自己的矩阵和对应像素，就得到 $W\bm X=\bm q$，其中 $W$ 有四行三列。若三台相机都看到这个点，叠起来就是六个方程，未知数仍只有 $X,Y,Z$。

\textbf{为什么用最小二乘？}即使有六个方程，理想观测也可以同时满足它们。实际像素定位有误差，反投影视线往往错开，讲义使用伪逆求代数残差平方和最小的点：
\[
 \hat{\bm X}=W^+\bm q,\qquad
 \hat{\bm X}\in\arg\min_{\bm X}\|W\bm X-\bm q\|_2^2.
\]
当 $W$ 满列秩时，解唯一，且 $W^+=(W^TW)^{-1}W^T$。两台相机若共用一个光心，只靠旋转拍摄，仍缺少确定深度所需的基线。

\textbf{从一个点到点云。}先固定相机，用标定物的已知三维点及其二维像点求各相机的 $P$。换成要重建的场景后，匹配各图中的二维像点，再对每组匹配分别求出三维坐标。与\relatedtopic{SC4061-L12-T01}{相机标定}相比，两步使用同一投影关系，但已知量与待求量互换。步骤见\relatedexercise{SC4061-L13-E03}{标定与重建的顺序}。讲义 p.~40 用同一组 69 张照片展示相机位姿、稀疏点云和稠密点云；pp.~41--42 另给出两组 COLMAP 稀疏重建结果。

\subsection{一分钟回顾}

水平立体校正把对应点搜索缩到同一行。白墙上许多位置都能得到很低的 SSD，遮挡处则缺少真正的对应点。找到可靠匹配后，每个视图为同一个三维点提供两个线性方程；实际观测有误差时，用最小二乘综合各视图的约束。
